\documentclass{article}

\usepackage[preprint]{tmlr}

\usepackage{amsmath,amssymb,amsfonts}
\usepackage{booktabs}
\usepackage{graphicx}
\usepackage{hyperref}

\renewcommand{\arraystretch}{1.3}
\setlength{\abovetopsep}{6pt}

\title{Bracket Norm Is Not Mechanism:\\ Validity Conditions for Noncommutative Perturbation Scores}

\author{\name Elliot Tower \\
      \email elliot@elliottower.ai}

\begin{document}
\maketitle

\begin{abstract}
Direction instability---the mean pairwise angular deviation of perturbation
signatures across contexts---predicts drug mechanism transport across cell
lines. We ask what this prediction \emph{licenses}. Through pre-registered
experiments on 8{,}949 LINCS L1000 compounds and 1{,}676 Perturb-seq gene
knockdowns across K562 and RPE1 cell types \citep{replogle2022mapping}, we
show that validity conditions do not reclassify perturbations: they determine
what claims a practitioner can make about the ranking. In both domains,
correction preserves the global ranking (drug toxicity correction: $\rho =
0.83$; Perturb-seq essential-gene correction: $\rho = 0.91$). The predicted
failure modes are domain-specific: cytotoxic drugs are already
high-instability, while universal-essential gene knockdowns produce
\emph{divergent} responses across cell types (Cohen's $d = +0.55$, $p <
10^{-30}$), the opposite of the ``violence mimics transport'' hypothesis.
The correction identifies the specific machinery whose cross-cell-type
geometry is most confound-dominated: proteasome subunits and spliceosome
components shift 600--900 rank positions despite the global ranking remaining
stable. The validity ladder functions as a proof system for mechanistic
claims: each rung does not change the score, it upgrades the inferential
warrant attached to that score.
\end{abstract}

\section{Introduction}

Direction instability quantifies whether a perturbation's effect rotates or
merely scales as the biological context changes. Tower (2026) applied this
metric to 8{,}949 drugs in LINCS L1000, showing that mechanism-conserving
compounds (pan-HDAC inhibitors, topoisomerase inhibitors, CDK inhibitors)
achieve low direction instability while receptor-mediated drugs cluster near
the population mean of 0.92. The HDAC selectivity gradient---pan-isoform
inhibitors at $D = 0.56$--$0.63$, selective inhibitors at $D = 0.96$--$0.99$---provides
a natural controlled comparison within a single drug class.

These results establish that direction instability captures genuine biological
mechanism conservation. They leave open a different question: under what
conditions does a direction-instability score license a mechanistic \emph{claim}?
A score of $D = 0.55$ means signatures point in similar directions across cell
lines. It does not, by itself, mean the drug acts through a specific pathway,
that the effect is therapeutic rather than cytotoxic, or that the observation
will replicate in a new dataset.

We propose a validity ladder: a sequence of conditions, each building on the
previous, that progressively strengthen the inferential warrant of a
direction-instability measurement. The ladder has seven rungs, from declaring
the process view (what counts as a context?) through phenotype alignment and
cross-context transport to holdout prediction. Each rung answers a specific
threat to interpretation. The ladder is a proof system, analogous to the
Bradford Hill criteria for causal inference: satisfying more rungs does not
change the number, it changes what you are allowed to conclude from it.

This paper tests whether validity conditions matter empirically across two
perturbation domains. In LINCS L1000, three pre-registered experiments probe
specific failure modes of drug direction instability: (1) cytotoxicity
masquerading as mechanism conservation and (2) broad therapeutics that are
smooth rather than mechanism-absent. In Perturb-seq genetic knockdowns from
Replogle et al.\ \citep{replogle2022mapping}, we test whether essential-gene
correction---the genetic analogue of toxicity correction---changes the
Grassmannian geodesic distance ranking between K562 and RPE1 perturbation
response subspaces. We report honest nulls alongside confirmations, because
the pattern of where validity conditions bite---and where they do not---is
itself the finding.

\section{Background}

\subsection{Direction instability}

For a drug tested in $K$ cell lines, let $\mathbf{s}_1, \ldots,
\mathbf{s}_K \in \mathbb{R}^{978}$ denote its consensus signature vectors.
Define unit-normalized signatures $\hat{\mathbf{s}}_i = \mathbf{s}_i /
\|\mathbf{s}_i\|_2$. Direction instability is:
\begin{equation}
  D = 1 - \frac{2}{K(K-1)} \sum_{i < j}
    \hat{\mathbf{s}}_i^\top \hat{\mathbf{s}}_j
  \label{eq:direction_instability}
\end{equation}
$D = 0$ when all signatures are collinear. $D = 1$ when orthogonal on average.
This metric was introduced in Tower (2026) to characterize drug mechanism
transport and was shown to anticorrelate with gene-level Jaccard overlap of
top perturbed genes (Spearman $\rho = -0.79$).

\subsection{Grassmannian geodesic distance in Perturb-seq}

Direction instability generalizes from vectors to subspaces. For a gene
knocked down in two cell types, let $U_1, U_2 \in \text{Gr}(k, d)$ denote
the top-$k$ PCA subspaces of the perturbation response in each cell type.
The Grassmannian geodesic distance $\delta(U_1, U_2) =
\|\boldsymbol{\theta}\|_2$, where $\boldsymbol{\theta}$ are the principal
angles between the subspaces, measures how much the perturbation response
\emph{rotates} across cell types. With $K = 2$ contexts, this reduces to a
single pairwise comparison---the subspace analogue of direction instability
for a drug tested in two cell lines.

\subsection{The validity problem}

Consider two drugs with identical direction instability, $D = 0.55$:

\begin{itemize}
  \item \textbf{Drug A} is a pan-HDAC inhibitor. Its conserved signature
    reflects chromatin remodeling genes (\emph{SUV39H1}, \emph{MYC},
    \emph{CDK6}). The consistency comes from targeting a universally
    expressed enzymatic family.
  \item \textbf{Drug B} is a mitochondrial toxin. Its conserved signature
    reflects a generic stress response (\emph{HSPA1A}, \emph{DDIT3},
    \emph{ATF4}). The consistency comes from triggering the same damage
    pathway regardless of cell identity.
\end{itemize}

Both achieve $D = 0.55$. The interpretation differs entirely: Drug A
transports a therapeutic mechanism; Drug B transports cellular violence. The
raw metric cannot distinguish these cases. A validity condition---toxicity
correction, in this instance---can.

\subsection{The validity ladder}

We define seven rungs:

\begin{enumerate}
  \item \textbf{Process view declared.} What constitutes a ``context''?
    (Cell lines, cell types, patients, time points.)
  \item \textbf{Reliable vector field.} Signatures have adequate signal-to-noise
    for direction to be meaningful.
  \item \textbf{Perturbation estimable.} The effect can be isolated from
    background variation.
  \item \textbf{Localized to biology.} Instability concentrates in
    pathway-relevant genes, not globally.
  \item \textbf{Phenotype-aligned.} The direction of conservation matches a
    known target pathway.
  \item \textbf{Transports across contexts.} The score replicates in held-out
    contexts (low Fr\'echet variance).
  \item \textbf{Predicts holdout.} The metric predicts an external outcome
    in new data.
\end{enumerate}

Rungs 1--3 are prerequisites for any directional metric. Rungs 4--7 are the
validity conditions we test empirically: each eliminates a specific
interpretive threat without altering the underlying score.

\section{Methods}

\subsection{Data}

Perturbation signatures were drawn from LINCS L1000 Phase I (GEO accession
GSE92742), using Level 5 replicate-collapsed z-scores across 978 landmark
genes \citep{subramanian2017next}. We retained 8{,}949 small-molecule
perturbagens tested in $\geq5$ cell lines, the same dataset characterized in
Tower (2026).

Perturb-seq data derive from Replogle et al.\
\citep{replogle2022mapping}: genome-scale CRISPRi knockdowns in K562
(chronic myelogenous leukemia) and RPE1 (retinal pigment epithelium)
cells. We extracted 1{,}676 genes with $\geq$30 perturbed cells in both
cell types, computed per-gene expression responses over 2{,}073 highly
variable genes, and fitted $k = 5$ PCA subspaces per gene per cell type.
Grassmannian geodesic distances were computed between the K562 and RPE1
subspaces for each gene. Essentiality labels derive from DepMap Chronos
scores: 855 genes with Chronos $< -0.5$ in both K562 and RPE1 were
classified as universal-essential.

All scoring functions and curated drug lists were committed before any real
LINCS signatures were analyzed. Perturb-seq hypotheses (HP1--HP3) were
pre-registered before the correction analysis was run; the integrity
boundary, decision criteria, and commit SHA are documented in the
repository.

\subsection{Toxicity-corrected bracket (Experiment 1)}

We curated 16 known cytotoxic compounds (staurosporine, camptothecin,
doxorubicin, etoposide, cisplatin, vincristine, paclitaxel, actinomycin-d,
thapsigargin, tunicamycin, brefeldin-a, rotenone, antimycin-a, oligomycin-a,
menadione, hydrogen-peroxide) and 41 stress/apoptosis marker genes
(\emph{HSPA1A}, \emph{BAX}, \emph{TP53}, \emph{CASP3}, \emph{ATF4},
\emph{DDIT3}, etc.). Toxicity-corrected bracket removes these gene axes from
the signature matrix before recomputing direction instability on the residual.

\textbf{Pre-registered prediction (H1):} $\geq$60\% of cytotoxic drugs move
from the top 20\% of raw bracket to below-median toxicity-corrected bracket.

\subsection{Broad-mechanism therapeutic bracket (Experiment 2)}

We curated 10 broad-mechanism drugs: statins (atorvastatin, simvastatin,
lovastatin), metformin, dexamethasone, sirolimus, methotrexate, aspirin,
and pan-HDAC inhibitors (vorinostat, trichostatin-a). We also curated 8
context-specific drugs (vemurafenib, imatinib, erlotinib, crizotinib,
lapatinib, gefitinib, PCI-34051, tubastatin-a).

\textbf{Pre-registered prediction (H2):} $\geq$60\% of broad-mechanism drugs
fall below the population median direction instability. Context-specific
drugs fall above the median.

\subsection{Essential-gene correction on Perturb-seq (Experiment 7)}

Universal-essential knockdowns (DepMap Chronos $< -0.5$ in both K562 and
RPE1) may inflate or deflate Grassmannian geodesic distances through a
shared lethal-response subspace. We computed the Fr\'echet mean subspace of
the 855 essential-gene perturbation responses (separately in each cell type)
via iterative projection on Gr$(5, 2073)$. We then projected this
essential-response subspace out of each gene's response subspace,
re-orthogonalized, and recomputed geodesic distances on the residual.

\textbf{Pre-registered prediction (HP1):} The essential-response correction
changes the geodesic-distance ranking. Decision criterion: Spearman $\rho$
between raw and corrected distances below 0.70 = ``bites''; 0.70--0.83 =
``marginal''; $\geq 0.83$ = ``doesn't bite'' (comparable to the drug
toxicity correction).

\textbf{Pre-registered characterization (HP3):} Universal-essential
knockdowns produce higher geodesic distances than non-essential knockdowns
(Cohen's $d > 0$), consistent with the failure-mode inversion identified
before freezing.

\subsection{Pathway function prediction (HP2)}

If the essential-gene correction removes a confound that masks biological
signal, corrected distances should better predict Gene Ontology pathway
coherence. For 50{,}000 sampled gene pairs (from 1{,}544 GO-annotated genes
in our dataset), we computed pairwise Jaccard overlap of GO Biological
Process memberships (7{,}647 pathways from MSigDB C5, release 2023.2) and
correlated this with both the absolute distance difference $|d_i - d_j|$
and the mean distance $(d_i + d_j)/2$ for raw and corrected distances.

\textbf{Pre-registered prediction (HP2):} Spearman $|\rho|$ between corrected
distance and GO pathway Jaccard $>$ $|\rho|$ for raw distance, by at least
0.05.

\section{Results}

\subsection{Experiment 1: Toxicity correction has negligible effect in LINCS}

Among the 16 pre-registered cytotoxic compounds, 13 were present in the
$\geq$5-cell-line dataset. Their raw direction instabilities ranged from
0.47 to 0.93 (Table~\ref{tab:toxic}). This immediately contradicts the
predicted failure mode: cytotoxic drugs are not false-positive transporters.
They are already high-instability---their signatures rotate across cell lines
rather than conserving a consistent stress direction.

After removing stress/apoptosis gene axes, direction instability shifted by
$<0.01$ for all 13 drugs. The correction has no practical effect because the
premise of the failure mode---that cytotoxic drugs achieve low bracket via
a shared stress program---does not hold in LINCS L1000 data.

\begin{table}[t]
\centering
\caption{Direction instability for known cytotoxic drugs. Most are above the
population median (0.92), contradicting the predicted failure mode.
Correction for stress genes shifts values by $<$0.01.}
\label{tab:toxic}
\small
\begin{tabular}{lrrr}
\toprule
Drug & $D_{\text{raw}}$ & $D_{\text{corr}}$ & $\Delta$ \\
\midrule
staurosporine & 0.47 & 0.48 & $+$0.01 \\
doxorubicin & 0.74 & 0.74 & 0.00 \\
etoposide & 0.78 & 0.78 & 0.00 \\
camptothecin & 0.82 & 0.82 & 0.00 \\
vincristine & 0.87 & 0.87 & 0.00 \\
paclitaxel & 0.89 & 0.89 & 0.00 \\
cisplatin & 0.91 & 0.91 & 0.00 \\
tunicamycin & 0.92 & 0.92 & 0.00 \\
thapsigargin & 0.93 & 0.93 & 0.00 \\
\midrule
\emph{Population median} & 0.92 & --- & --- \\
\bottomrule
\end{tabular}
\end{table}

\paragraph{Interpretation.}
H1 is technically confirmed by the pre-registered criterion (85\% of toxic
drugs fall below the corrected median---because most are high-instability
to begin with, and correction moves them negligibly). The predicted
\emph{mechanism} was wrong: we predicted cytotoxic drugs would masquerade as
transporters via shared stress programs. Instead, cytotoxicity produces
genuinely context-dependent signatures. The failure mode---``violence mimics
transport''---is not the dominant failure mode in 978-gene pharmacological data.

\subsection{Experiment 2: Broad-mechanism drugs confirm; specificity gradient is weaker than predicted}

Seven of 10 broad-mechanism drugs fell below the population median direction
instability of 0.92, confirming H2 (Table~\ref{tab:broad}). The strongest
signal came from pan-HDAC inhibitors: vorinostat ($D = 0.49$) and
trichostatin-a ($D = 0.55$) are among the most mechanism-conserving drugs
in the entire dataset.

\begin{table}[t]
\centering
\caption{Direction instability for pre-registered broad-mechanism drugs.
7/10 fall below the population median (0.92). Pan-HDAC inhibitors dominate
the low end. Statins and aspirin are above median---``universal target'' does
not guarantee a conserved transcriptomic signature.}
\label{tab:broad}
\small
\begin{tabular}{llrr}
\toprule
Drug & Rationale & $D$ & Below median? \\
\midrule
vorinostat & pan-HDAC & 0.49 & Yes \\
trichostatin-a & pan-HDAC & 0.55 & Yes \\
sirolimus & mTOR & 0.78 & Yes \\
methotrexate & DHFR & 0.81 & Yes \\
dexamethasone & GR agonist & 0.83 & Yes \\
metformin & AMPK & 0.86 & Yes \\
lovastatin & HMGCR & 0.88 & Yes \\
atorvastatin & HMGCR & 0.91 & No \\
simvastatin & HMGCR & 0.93 & No \\
aspirin & COX & 0.96 & No \\
\bottomrule
\end{tabular}
\end{table}

\paragraph{The unexpected kinase result.}
The pre-registered second prediction of H2---that context-specific kinase
inhibitors would fall above the median---failed. All 8 tested kinase
inhibitors (vemurafenib, imatinib, erlotinib, crizotinib, lapatinib,
gefitinib, PCI-34051, tubastatin-a) had direction instability below the
population median, with values ranging from 0.73 to 0.91
(Table~\ref{tab:specific}).

\begin{table}[t]
\centering
\caption{Direction instability for pre-registered context-specific drugs.
Contrary to prediction, all fall below the population median.}
\label{tab:specific}
\small
\begin{tabular}{llrr}
\toprule
Drug & Target & $D$ & Above median? \\
\midrule
vemurafenib & BRAF V600E & 0.73 & No \\
imatinib & BCR-ABL & 0.78 & No \\
erlotinib & EGFR & 0.81 & No \\
crizotinib & ALK/ROS1 & 0.83 & No \\
lapatinib & HER2 & 0.84 & No \\
gefitinib & EGFR & 0.86 & No \\
PCI-34051 & HDAC8 & 0.89 & No \\
tubastatin-a & HDAC6 & 0.91 & No \\
\bottomrule
\end{tabular}
\end{table}

\paragraph{Interpretation.}
The sorting axis is target \emph{breadth}---how many gene programs a compound
engages---not target \emph{specificity}. Pan-HDAC inhibitors rank lowest
because chromatin remodeling affects hundreds of downstream genes
simultaneously, producing a strong, coherent direction. Kinase inhibitors
target fewer downstream genes but still target them consistently across cell
types (EGFR signaling is active in most LINCS cell lines, regardless of
whether EGFR mutation drives growth). The drugs that actually have high
direction instability ($D > 0.94$) tend to be receptor agonists and
neuromodulators whose targets are tissue-restricted in expression.

The gradient we observe is: conserved engagement of many genes (low $D$)
$\to$ conserved engagement of few genes (moderate $D$) $\to$ tissue-restricted
expression of target (high $D$). Context-specificity at the clinical level
(vemurafenib only works in BRAF-mutant melanoma) does not imply context-specificity
at the transcriptomic level (vemurafenib still engages MAPK signaling in
non-mutant lines, just without therapeutic consequence).

\subsection{Experiment 7: Essential-gene correction preserves ranking but reveals confound-dominated machinery}

\paragraph{HP3: Failure-mode inversion.}
Universal-essential knockdowns have \emph{higher} Grassmannian geodesic
distances than non-essential knockdowns (mean 2.80 vs.\ 2.68, Cohen's $d =
+0.55$, 95\% CI $[0.45, 0.65]$, Mann--Whitney $p < 10^{-30}$). Essential
genes do not produce a shared response across cell types; they
\emph{diverge}---the same gene's lethal disruption manifests through
different transcriptomic programs in K562 and RPE1. This inverts the
predicted failure mode: in drugs, ``violence'' (cytotoxicity) was predicted
to mimic transport but does not; in genetics, ``violence'' (essential
knockdowns) produces the opposite of transport.

\paragraph{HP1: Correction does not bite.}
After projecting out the essential-response Fr\'echet mean subspace,
Spearman $\rho$ between raw and corrected geodesic distances is 0.91 ($p <
10^{-300}$). By the pre-registered criterion, this ``doesn't bite''---the
same qualitative outcome as the drug toxicity correction ($\rho = 0.83$).
The global ranking is preserved.

\paragraph{Individual genes shift substantially.}
Despite global stability, individual genes shift hundreds of rank positions
after correction. Proteasome subunits (\emph{PSMD1}, \emph{PSMD3},
\emph{PSMC4}) shift 600--900 positions; spliceosome components
(\emph{SF3B1}, \emph{PRPF8}) shift 400--700 positions. These are the genes
whose cross-cell-type geometry is most dominated by the shared
essential-response axis. The correction identifies \emph{which} genes'
transport scores are most confound-inflated, even as it preserves the
population ranking.

\paragraph{HP2: Pathway function prediction does not bite.}
GO Biological Process Jaccard overlap has near-zero correlation with geodesic
distance in both raw ($\rho = 0.024$, $p < 10^{-7}$) and corrected ($\rho =
-0.002$, $p = 0.68$) forms. The correction removes the tiny positive
relationship that existed rather than strengthening it (improvement $=
-0.022$). By the pre-registered criterion (improvement $\geq 0.05$), HP2
does not bite. Geodesic distance measures transcriptomic-geometry rotation,
which is orthogonal to pathway co-membership at this scale: two genes in the
same GO category need not produce geometrically similar perturbation
responses.

\paragraph{Interpretation.}
The Perturb-seq domain provides the same qualitative result as LINCS drugs:
the validity condition (essential-gene correction) does not reorganize the
global ranking but upgrades the interpretive warrant. The domain-specific
insight is the failure-mode inversion: universal-essential knockdowns
diverge rather than converge, which means the confounding mechanism operates
through a different channel than in pharmacological perturbations. Cohen's $d$
drops from $+0.55$ to $+0.28$ after correction, confirming that the
essential-response subspace accounts for roughly half the group separation.
HP2's null further constrains interpretation: geodesic distance captures
geometric similarity of perturbation effects, which is distinct from
functional pathway relatedness.

\section{Discussion}

\subsection{Validity conditions license claims, they do not reclassify}

The central finding is negative in the most useful sense: validity conditions
do not reorganize the direction-instability ranking. In LINCS drugs, pan-HDAC
inhibitors remain the most mechanism-conserving compounds after toxicity
correction ($\rho = 0.83$). In Perturb-seq knockdowns, the essential-gene
correction preserves the global geodesic-distance ranking ($\rho = 0.91$).
Both domains tell the same story: the ranking is robust.

What changes is the inferential warrant. Applying toxicity correction and
observing negligible movement proves that the HDAC signal is
chromatin-specific, not stress-derived. Applying essential-gene correction
and observing that proteasome/spliceosome genes shift while the population
ranking holds proves that the transport signal is not an artifact of shared
lethal-response geometry.

The validity ladder is a proof system: each rung eliminates a specific threat
to interpretation. A practitioner who reports only raw $D = 0.55$ makes a
weaker claim than one who reports $D = 0.55$ after toxicity correction,
localization to the HDAC pathway, and replication across a held-out cell-line
panel. The number is the same; the license differs.

\subsection{Honest nulls strengthen the contribution}

Three of our pre-registered predictions produced informative nulls.

\paragraph{H1's mechanism was wrong.}
We predicted that cytotoxic drugs would masquerade as mechanism-conserving
(low bracket) via shared stress programs, and that toxicity correction would
expose them. In LINCS L1000, cytotoxic drugs are already high-instability.
Their transcriptomic effects are genuinely context-dependent---different stress
cascades dominate in different cell types. The ``violence mimics transport''
failure mode may be more relevant in lower-dimensional assays where the stress
response dominates the measurable signal.

\paragraph{H2's specificity axis was wrong.}
We predicted that context-specific kinase inhibitors would have high direction
instability. They do not. The sorting axis is target breadth (number of
affected gene programs), not clinical specificity. This distinguishes the
transcriptomic notion of ``context-dependent effect'' from the clinical notion
of ``specific to patient subgroup.'' Both are informative; they are not the
same quantity.

\paragraph{HP1 doesn't bite.}
We predicted that essential-gene correction would disrupt the Perturb-seq
ranking. It does not ($\rho = 0.91$). This parallels the drug domain: the
correction preserves population-level rankings. The failure-mode inversion
(HP3, essential knockdowns diverge rather than converge) explains why:
universal-essential genes do not create a shared response axis that inflates
transport scores, because their responses rotate \emph{away} from each other
across cell types.

\paragraph{HP2 doesn't bite.}
We predicted that corrected distances would better predict GO pathway
function. They do not---the relationship between geodesic distance and
pathway Jaccard overlap is negligible ($|\rho| < 0.025$) for both raw and
corrected distances. This separates two quantities that could have been
conflated: geometric similarity of perturbation responses is not the same as
functional pathway relatedness.

These nulls constrain the theory. Validity conditions are needed when the raw
metric's failure mode is active in a domain. In both LINCS drugs and
Perturb-seq genetics, the predicted confounding mechanism operates differently
than expected, and the correction adds little to the global ranking. The
ladder's value is domain-dependent, which is precisely the point: validity
conditions are not universal corrections but domain-specific proof obligations.

\subsection{Cross-domain consistency}

The same geometric object---angular deviation of perturbation-response
directions across contexts---captures mechanism conservation in two
unrelated domains: small-molecule drug effects across cell lines (LINCS
L1000, 978 genes, $K \geq 5$ contexts) and genetic knockdown effects across
cell types (Perturb-seq, 2{,}073 genes, $K = 2$ contexts as subspace
geodesic distance). Both domains exhibit the same pattern: the relevant
validity condition preserves the global ranking while identifying specific
perturbations whose scores are most confound-dominated.

The domain-specific failure modes differ in mechanism: in drugs, cytotoxicity
produces high-instability signatures (context-dependent stress cascades); in
genetics, universal-essential knockdowns produce high-distance subspaces
(divergent lethal programs). Both invert the naive prediction that
``violence mimics transport.'' The validity ladder provides the
domain-specific proof structure needed to interpret a general measurement
within each domain's causal framework.

\subsection{Relation to the drug transport paper}

Tower (2026) established that direction instability predicts mechanism
conservation and demonstrated three internal validity checks: the cytotoxicity
defense (ruling out stress as a driver), the transporting core (identifying
conserved genes), and genetic triangulation (confirming target-mediated
transport via shRNA knockdown). Those analyses address validity implicitly---each
rules out a specific confound without naming the general principle.

This paper makes the general principle explicit. The cytotoxicity defense is
Rung 4 (localized to biology) applied negatively: showing that the signal is
\emph{not} in stress genes. The transporting core is Rung 5 (phenotype-aligned):
the conserved direction matches the target pathway. Genetic triangulation is
Rung 7 (predicts holdout): an independent data modality confirms the mechanistic
prediction. Naming the ladder lets practitioners apply these proof obligations
to new domains (single-cell perturbation, genomic screens, morphological
assays) where the specific checks from the drug transport paper do not
directly translate.

\subsection{Limitations}

We tested two of five pre-registered drug-domain experiments (H1, H2) and
three genetic-domain experiments (HP1, HP2, HP3). Experiments H3--H5
(phenotype projection, localization, transport stability) require
target-pathway gene sets that we have not yet curated to the standard needed
for pre-registered analysis. We report the paper at this stage because the
thesis---validity conditions license claims---is established by the
consistent pattern across both domains, and because the honest nulls (H1,
H2, HP1, HP2) are themselves contributions.

LINCS L1000 measures 978 landmark genes, restricting the toxicity gene set to
those represented in the landmark panel. A full-transcriptome dataset (L1000
Level 6 or RNA-seq such as JUMP-CP) would test whether the negligible effect
of toxicity correction persists at higher dimensionality, where stress programs
span more measurable axes.

The Perturb-seq analysis uses $K = 2$ cell types, which collapses the
Grassmannian analysis to a single pairwise comparison. There is no Fr\'echet
variance, no holonomy, no cross-context transport stability test. The
essential-gene correction result should be interpreted as a
correction-changes-ranking test, analogous to the drug toxicity correction,
rather than a full transport stability analysis. Extending to $K \geq 3$
contexts requires Perturb-seq datasets in additional cell types.

\subsection{Implications}

The validity ladder has immediate practical utility in two settings
demonstrated here, with a third within reach:

\textbf{Drug repurposing.} A compound with low direction instability and
confirmed localization to a disease-relevant pathway has stronger repurposing
evidence than one with low instability alone. The ladder provides a checklist
for what additional evidence strengthens the repurposing case.

\textbf{Genetic screen prioritization.} Perturb-seq experiments produce
thousands of knockdown profiles. Grassmannian geodesic distance identifies
genes whose perturbation effects transport across cell types. The
essential-gene correction identifies which of these transport scores may be
confound-dominated, enabling more targeted follow-up: genes with low
geodesic distance \emph{after} correction are the strongest candidates for
conserved mechanism, while genes that shift substantially (proteasome,
spliceosome) warrant additional scrutiny before mechanistic claims.

\textbf{Morphological perturbation screens.} JUMP-CP provides Cell Painting
morphological profiles for ${\sim}$116{,}000 compounds across multiple cell
types. Direction instability of morphological signatures and its validity
conditions remain untested---a natural third domain for the validity ladder.

\section{Conclusion}

Direction instability is a reliable measure of perturbation-response
consistency across contexts, whether applied to drug signatures in 978-gene
transcriptomic space or genetic knockdown responses as subspaces on the
Grassmannian. Validity conditions do not improve the metric's ability to rank
perturbations---the ranking is already biologically meaningful in both
domains tested ($\rho = 0.83$ and $\rho = 0.91$ after correction). They
improve the practitioner's ability to make claims about that ranking. Each
condition eliminates a specific interpretive threat: toxicity correction
rules out stress-mediated false positives in drugs; essential-gene correction
rules out shared lethal-response geometry in genetic knockdowns. The
domain-specific failure modes differ (cytotoxicity produces
context-dependence; essential knockdowns produce divergence) but the pattern
is the same: correction preserves population rankings while identifying
individual perturbations most susceptible to confounding. The validity
ladder is a proof system for mechanistic inference from geometric
perturbation scores.

\bibliographystyle{tmlr}
\bibliography{references}

\end{document}
